\section{Version A: the material intensity of capital tied to the final good}\label{sec:sp:A}

Section~\ref{sec:sp:setup} reduced the process-level accounting to a single physical relation, the
ledger~\eqref{eq:sp:setup:ledger}, and left one modelling choice open: how the material intensity of
new capital goods, $\Phi^I$, is determined. This section develops the planner's problem in full under
Version A, in which capital goods inherit the material intensity of final-good activity in general,
$\Phi^I=\Omega\phi^I$. Section~\ref{sec:sp:B} does the same under Version B, in which $\Phi^I$ is
exogenous, and Section~\ref{sec:sp:B:compare} compares the two.


\subsection{The closure and what it determines}\label{sec:sp:A:closure}

Under Version A the material content of capital goods per unit of expenditure is not pinned down
independently of the material content of final-good activity in general. Capital goods are made from
the same composite good as everything else, so
\begin{align}
\Phi^I &= \Omega\phi^I,
\label{eq:sp:A:closure}
\end{align}
with $\phi^I$ an exogenous constant: capital-goods production is $\phi^I$ times as material-intensive
as final-good activity on average, $\phi^I=1$ being the parsimonious case. The level $\Omega$ is then
determined by the aggregate balance~\eqref{eq:sp:setup:ybalance_OmegaPhi},
\begin{align}
R &= \Omega\left(\mathcal D+\phi^IG\right),
&
\mathcal D &= \omega^YY+\omega^{N}C^N+\omega^{D}C^D+\omega^WC^W+\omega^CC,
\label{eq:sp:A:Omega}
\end{align}
where $\mathcal D$ gathers every flow measured in final-good units and therefore carrying embodied
material drawn through $R$. Three consequences of this closure shape the problem below, and each is
worth having in view before the algebra.

\smalltitle{(i) The non-negativity restriction is vacuous.} Solving~\eqref{eq:sp:A:Omega} gives
$\Omega=R/\left(\mathcal D+\phi^IG\right)\ge0$ automatically, since $R>0$ and both terms in the
denominator are non-negative. The intensity of capital goods scales down to accommodate whatever
gross investment the planner chooses, rather than the planner having to find enough material to meet
a fixed per-unit requirement. The constraint $\Phi^IG\le R$, which under Version B is a genuine
restriction on the feasible set, can therefore be dropped from the problem entirely.

\smalltitle{(ii) The ledger is no longer linear in $W$.} Because $\Omega$ depends on $W$ --- through
$\mathcal D$, which contains $\omega^WC^W=\omega^Wc^W\left(\varpi\right)W$, and through $R=N+a\varpi W$
--- so does $\Phi^I$, and hence so does $\dot M^K$. The collected
ledger~\eqref{eq:sp:setup:ledger-collected} does not solve in one line; clearing denominators turns it
into the quadratic~\eqref{eq:sp:setup:ledger-quadratic}, which has a positive leading coefficient and a
negative constant term and therefore exactly one positive root. $W$ is still determined uniquely and
in closed form, but conservation of mass no longer fixes it by a division.

\smalltitle{(iii) The composition of final-good spending is a margin.} This is the substantive
content of the closure and the reason the section is worth writing separately. Spending one more unit
of the final good on any activity $j$ raises $\mathcal D$ by $\omega^j$, which
by~\eqref{eq:sp:A:Omega} \emph{dilutes} $\Omega$ for given $R$, which lowers $\Phi^I$, which means
less of the incoming mass is embodied in new capital and more of it is diffused as waste today.
Total mass is conserved throughout --- what the composition of spending changes is the \emph{split}
of that mass between the durable stock and the current waste flow. That is a real margin, and the
optimality conditions below price it.


\subsection{The planner's problem}\label{sec:sp:A:problem}

\smalltitle{Controls and states.} The planner chooses paths for the eight controls
\begin{align}
\left\{C,\;I,\;N,\;D,\;\varpi,\;K^R,\;W,\;\Omega\right\},
\label{eq:sp:A:controls}
\end{align}
which determine the remaining flows through the definitions
\begin{align}
R^R &= \mathcal R\left(\varpi W,K^R\right),
&
R &= N+R^R,
&
Y &= F\left(K-K^R,\,R,\,P\right),
&
G &= I+\delta K,
\label{eq:sp:A:definitions}
\end{align}
and it carries the five states
\begin{align}
\left\{K,\;S,\;X,\;P,\;M^K\right\},
\label{eq:sp:A:states}
\end{align}
with $K_0,S_0,X_0,P_0,M^K_0$ given.

\smalltitle{Why $W$ and $\Omega$ are listed as controls.} Neither is an independent object: the
ledger determines $W$ residually and~\eqref{eq:sp:A:Omega} determines $\Omega$ residually. Listing
them among the controls and imposing the two relations as equality constraints is the equivalent
formulation that keeps their shadow prices visible. Those prices are what the section is about ---
$p^W$ is the object every environmental instrument in the market economy is built from, and the
multiplier on~\eqref{eq:sp:A:Omega} is what the entire difference between the two versions turns out
to be expressible in terms of. Remark~\ref{rem:sp:A:reduced} states the equivalence with the reduced
formulation.

\begin{problem}[Planner, Version A]\label{prob:sp:A}
Choose $\left\{C,I,N,D,\varpi,K^R,W,\Omega\right\}_{t\ge0}$ to maximise
$U_0=\int_0^\infty\left[u(C)-v(P)\right]e^{-\rho t}dt$ subject to, at every $t$,
\begin{align}
&\text{\emph{goods:}}
&
F\left(K-K^R,R,P\right) &= C+I+\delta K+C^N(N,S)+C^D(D,X)+c^W\!\left(\varpi\right)W,
\label{eq:sp:A:goods}
\\[0.2em]
&\text{\emph{ledger:}}
&
W &= R-\Omega\phi^I G+\delta M^K+\mathcal N,
\label{eq:sp:A:ledger}
\\[0.2em]
&\text{\emph{intensity:}}
&
R &= \Omega\left(\mathcal D+\phi^IG\right),
\label{eq:sp:A:intensity}
\end{align}
and to the laws of motion
\begin{align}
\dot K &= I,
&
\dot S &= D-N,
&
\dot X &= D,
&
\dot P &= \Xi-\theta(P)P,
&
\dot M^K &= \Omega\phi^I G-\delta M^K,
\label{eq:sp:A:states-motion}
\end{align}
with $\mathcal N=\Omega^{N,S}N+\Omega^{D,S}D$, $c^W(\varpi)=c^c+c^T(\varpi)$, $\mathcal D$ as
in~\eqref{eq:sp:A:Omega}, the definitions in~\eqref{eq:sp:A:definitions}, the emission flow
\begin{align}
\Xi &= \left(1-\varpi+d^W\varpi\right)W-d^WR^R,
\label{eq:sp:A:Xi}
\end{align}
and $C,I+\delta K,N,D,K^R,W,\Omega\ge0$, $\varpi\in\left[0,1\right]$.
\end{problem}

Three features of this statement deserve marking before any algebra.

First, \eqref{eq:sp:A:ledger} is an \emph{identity}, not a technology. It says where the mass went,
and it holds along every conceivable path; what makes it bite is that it ties the waste flow $W$ ---
which is costly to collect, costly to treat and damaging when released --- to the material flow $R$
that production requires. Nothing in it can be relaxed by choosing differently; it can only be made
less costly by choosing a path on which less mass has to be moved.

Second, \eqref{eq:sp:A:intensity} is likewise an accounting definition rather than a restriction,
as~(i) above records: it determines $\Omega$ and never binds. Its multiplier is nonetheless nonzero,
and this is not a contradiction. The multiplier is not the price of a scarce resource; it is the price
of the \emph{feedback} the definition creates, namely that the composition of spending changes
$\Omega$ and hence changes how much of the incoming mass is stored rather than shed.

Third, $\Xi$ is written in the form~\eqref{eq:sp:A:Xi}, linear in $W$ and $R^R$ separately, so that no
derivative of $a(\cdot)$ ever appears inside a pollution term. All the curvature of the recycling
technology enters through $\mathcal R$ and through $\alpha$ and $a'$ alone; recall
from~\eqref{eq:sp:setup:recycling-derivatives} that $\mathcal R_T=\alpha\equiv a-a'x$ is the marginal
recycling yield and $\mathcal R_{K^R}=a'$.

\smalltitle{The Lagrangian.} Write the current-value Hamiltonian with the static constraints
adjoined. It is convenient to normalise every multiplier by $\lambda$, the multiplier on the goods
constraint~\eqref{eq:sp:A:goods}, so that all shadow prices below are measured in units of the final
good rather than in utils. Accordingly, let the current-value costates be
\begin{align}
\underbrace{\lambda q}_{K},
&&
\underbrace{\lambda p^S}_{S},
&&
\underbrace{\lambda p^X}_{X},
&&
\underbrace{-\lambda p^P}_{P},
&&
\underbrace{-\lambda p^M}_{M^K},
\label{eq:sp:A:costates-def}
\end{align}
the two minus signs anticipating that both $P$ and $M^K$ are liabilities, so that $p^P$ and $p^M$
read as shadow \emph{costs}; the signs are verified in Section~\ref{sec:sp:A:reading}, not assumed.
Let $\lambda p^W$ be the multiplier on the ledger~\eqref{eq:sp:A:ledger} written with $W$ on the left,
and $\lambda\zeta$ the multiplier on the intensity definition~\eqref{eq:sp:A:intensity} written with
$R$ on the left. Then
\begin{align}
\mathcal L
={}& u(C)-v(P)
\notag\\
&+\lambda\Big[F\left(K-K^R,R,P\right)-C-I-\delta K-C^N(N,S)-C^D(D,X)-c^W\!\left(\varpi\right)W\Big]
\notag\\
&+\lambda q\,I
\;+\;\lambda p^S\left(D-N\right)
\;+\;\lambda p^X D
\;-\;\lambda p^P\left[\Xi-\theta(P)P\right]
\;-\;\lambda p^M\left[\Omega\phi^IG-\delta M^K\right]
\notag\\
&+\lambda p^W\Big[W-N-R^R+\Omega\phi^IG-\delta M^K-\mathcal N\Big]
\;+\;\lambda \zeta\Big[N+R^R-\Omega\left(\mathcal D+\phi^IG\right)\Big],
\label{eq:sp:A:lagrangian}
\end{align}
with $R=N+R^R$, $R^R=\mathcal R\left(\varpi W,K^R\right)$, $G=I+\delta K$, $Y=F\left(K-K^R,R,P\right)$
and $C^W=c^W\!\left(\varpi\right)W$ substituted throughout --- the last two matter, because $\mathcal D$
contains both $\omega^YY$ and $\omega^WC^W$, so every control that moves output or the waste flow also
moves the waste base.

The sign convention on $p^W$ is worth stating explicitly, because it is the price everything else is
written in terms of. Perturb the ledger to $W=R-\Omega\phi^IG+\delta M^K+\mathcal N-\kappa$, that is,
suppose $\kappa$ tonnes of waste could be made to disappear at no cost. Then
$\partial\mathcal L^*/\partial\kappa=\lambda p^W$, so $p^W>0$ says exactly that waste is socially
costly, and $p^W$ measures how costly, per tonne, in final-good units.


\subsection{First-order conditions}\label{sec:sp:A:foc}

\smalltitle{The intensity margin, and the wedge.} Take the condition for $\Omega$ first, since it
determines $\zeta$ in closed form and everything else is then read off by substitution.
Differentiating~\eqref{eq:sp:A:lagrangian} with respect to $\Omega$ and dividing by $\lambda$,
\begin{align}
\phi^IG\left(p^W-p^M\right) &= \zeta\left(\mathcal D+\phi^IG\right)
\;=\;\zeta\,\frac{R}{\Omega},
\label{eq:sp:A:foc:Omega}
\end{align}
the second equality using~\eqref{eq:sp:A:intensity}. Define the share of material input that is
embodied in new capital,
\begin{align}
\chi &\equiv \frac{\Phi^IG}{R}\;=\;\frac{\Omega\phi^IG}{R}\;\in\;\left[0,1\right),
\qquad\text{so that}\qquad
\zeta = \chi\left(p^W-p^M\right),
\qquad
\Delta \equiv \zeta\Omega = \Omega\chi\left(p^W-p^M\right).
\label{eq:sp:A:zeta}
\end{align}
$\Delta$ is the object that does all the work below, and it is dimensionless: $\Omega$ is tonnes per
unit of final good and $p^W-p^M$ is final good per tonne. Its interpretation is the one given
in~(iii) of Section~\ref{sec:sp:A:closure}. Raising spending on activity $j$ by one unit dilutes
$\Omega$, shifts mass out of the durable stock and into today's waste flow, and the social value of
that shift is the difference between shedding a tonne now and shedding it later, $p^W-p^M$, scaled by
how much of the throughput is at stake, $\chi$. Note that $\Delta$ carries the sign of $p^W-p^M$, a
fact used repeatedly below.

The remaining conditions follow by differentiating~\eqref{eq:sp:A:lagrangian} with respect to each of
the other seven controls and dividing by $\lambda$. Two composite objects recur often enough to be
named at the outset:
\begin{align}
B &\equiv \left(1-\Delta\omega^Y\right)F_R-p^W+d^Wp^P+\zeta,
&
p^N &\equiv \left(1+\Delta\omega^N\right)C^N_N+p^S+\Omega^{N,S}p^W .
\label{eq:sp:A:Bvalue}
\end{align}
$B$ is the net social value of one tonne of \emph{secondary} material: what it produces
--- valued net of the waste base its output adds --- less the waste it will itself become, plus the
emission its recovery avoids, plus the dilution of $\Omega$ that the extra material permits.
$p^N$ is the full marginal social cost of one tonne of \emph{virgin} material delivered
to production: the resources spent extracting it, marked up because extraction effort is itself
final-good spending that adds to the waste base, plus the rent on the reserve it depletes, plus the
waste cost of the overburden its extraction displaces.

\smalltitle{Consumption.}
\begin{align}
u'\!\left(C\right) &= \lambda\left(1+\Delta\omega^C\right).
\label{eq:sp:A:foc:C}
\end{align}
Consumption carries a wedge. A unit consumed is a unit of final-good spending that adds $\omega^C$ to
the waste base, dilutes $\Omega$, and thereby diverts mass from the capital stock into the waste flow.

\smalltitle{Investment.}
\begin{align}
q &= 1-\Phi^I\left(1-\chi\right)\left(p^W-p^M\right).
\label{eq:sp:A:foc:I}
\end{align}
Building a unit of capital costs one unit of the final good and takes $\Phi^I$ tonnes out of the
current waste stream, worth $\Phi^Ip^W$; those tonnes return as waste over the life of the asset, at a
present value of $\Phi^Ip^M$. Capital is a materials sink, but a temporary one. The factor
$\left(1-\chi\right)$ damps the effect: raising gross formation embodies more mass in capital, but it
also raises $\phi^IG$ in the denominator of~\eqref{eq:sp:A:Omega}, diluting $\Omega$ and hence
$\Phi^I$ itself, and $\chi$ is exactly the fraction of the direct effect that this feedback offsets.

\smalltitle{Extraction.}
\begin{align}
\left(1-\Delta\omega^Y\right)F_R+\zeta-p^W &= \left(1+\Delta\omega^N\right)C^N_N+p^S+\Omega^{N,S}p^W,
\qquad\text{equivalently}\qquad
B &= p^N+d^Wp^P .
\label{eq:sp:A:foc:N}
\end{align}

\smalltitle{Exploration.}
\begin{align}
p^S+p^X &= \left(1+\Delta\omega^D\right)C^D_D+\Omega^{D,S}p^W .
\label{eq:sp:A:foc:D}
\end{align}

\smalltitle{Treatment share.}
\begin{align}
\alpha B+\left(1-d^W\right)p^P &= \left(1+\Delta\omega^W\right)\frac{dc^T}{d\varpi},
\label{eq:sp:A:foc:varpi}
\end{align}
for interior $\varpi\in\left(0,1\right)$, with the usual inequalities at the corners.

\smalltitle{Recycling capital.}
\begin{align}
\left(1-\Delta\omega^Y\right)F_K &= a'\!\left(x\right)B.
\label{eq:sp:A:foc:KR}
\end{align}

\smalltitle{Waste.}
\begin{align}
p^W &= \left(1+\Delta\omega^W\right)c^W\!\left(\varpi\right)
+\left(1-\varpi+d^W\varpi\right)p^P-\varpi\alpha B.
\label{eq:sp:A:foc:W}
\end{align}

\smalltitle{The pattern.} The seven conditions have a single organising principle, and it is worth
stating once rather than re-reading it off each equation. \emph{A unit of the final good spent on
activity $j$ has effective social cost $\left(1+\Delta\omega^j\right)$; a unit produced has effective
social value $\left(1-\Delta\omega^Y\right)$.} The asymmetry in sign is not an error: spending adds
to the waste base and dilutes $\Omega$, whereas producing output also adds $\omega^YY$ to the base,
so output is worth less than its marginal product rather than more. Setting $\Delta=\zeta=0$ removes
every wedge and returns the conditions of Version B.

Two of the conditions are worth expanding on immediately, because their form is not obvious.

\eqref{eq:sp:A:foc:W} defines $p^W$ implicitly, since $B$ contains $p^W$. Solving out
using~\eqref{eq:sp:A:foc:N} to replace $B$, the shadow cost of a tonne of gross waste
satisfies
\begin{align}
p^W\left(1+\Omega^{N,S}\varpi\alpha\right)
&= \left(1+\Delta\omega^W\right)c^W\!\left(\varpi\right)+\left(1-\varpi+d^W\varpi\right)p^P
-\varpi\alpha\Big[\left(1+\Delta\omega^N\right)C^N_N+p^S+d^Wp^P\Big].
\label{eq:sp:A:foc:W-solved}
\end{align}
A tonne of waste costs what it costs to collect and treat, plus the damage from the part of it that
reaches the environment, minus the virgin extraction it displaces at the margin. Unlike its Version~B
counterpart, this expression does involve the relative intensities, through $\Delta$; and since
$\Delta$ is itself proportional to $p^W-p^M$, \eqref{eq:sp:A:foc:W-solved} is a genuinely implicit
determination of $p^W$ rather than a formula for it.

\eqref{eq:sp:A:foc:varpi} is the treatment margin. Raising $\varpi$ by one point costs
$\left(1+\Delta\omega^W\right)dc^T/d\varpi$ per tonne of waste and buys two things: $\alpha$ tonnes of
secondary material worth $B$ each, and a reduction of the emitted fraction from $1$ to
$d^W$ on the tonne now treated. Since $c^T(0)=\left.dc^T/d\varpi\right|_{0}=0$
by~\eqref{eq:sp:setup:waste-cost}, the right-hand side vanishes at $\varpi=0$ while the left-hand side
is strictly positive whenever $p^P>0$; so \emph{some} treatment is always optimal, and the $\varpi=0$
corner is never reached. The upper corner $\varpi=1$ is reached iff
$\alpha B+\left(1-d^W\right)p^P\ge\left(1+\Delta\omega^W\right)\left.dc^T/d\varpi\right|_{1}$.


\subsection{Costate equations}\label{sec:sp:A:costates}

Each costate obeys $\dot\nu=\rho\nu-\partial\mathcal L/\partial(\text{state})$ with $\nu$ the
current-value costate in~\eqref{eq:sp:A:costates-def}. Because the costates carry the factor
$\lambda$, each such equation produces a term $\dot\lambda/\lambda$, and it is worth defining the
resulting discount rate once:
\begin{align}
r &\equiv \rho-\frac{\dot\lambda}{\lambda}.
\label{eq:sp:A:r}
\end{align}
Every goods-denominated shadow price below discounts at this same $r$, because every one of them has
been normalised by $\lambda$. What does \emph{not} hold under Version A is the further identification
of $r$ with the consumption rate of interest: by~\eqref{eq:sp:A:foc:C}, $\lambda\ne u'(C)$, and the
gap between them is the consumption wedge. This is taken up
after~\eqref{eq:sp:A:euler}.

\smalltitle{Capital.} Using
$\partial\mathcal L/\partial K=\lambda\left[\left(1-\Delta\omega^Y\right)F_K-\delta
-\delta\Phi^I\left(p^M-p^W+\zeta\right)\right]$ and~\eqref{eq:sp:A:foc:I} to collapse the bracket to
$\left(1-\Delta\omega^Y\right)F_K-\delta q$,
\begin{align}
r &= \frac{\left(1-\Delta\omega^Y\right)F_K+\dot q}{q}-\delta .
\label{eq:sp:A:costate:K}
\end{align}
Depreciation is charged at $q$, not at $1$: what wears out is a unit of installed capital, whose
social value is $q$, not a unit of the final good.

\smalltitle{Reserves.} With
$\partial\mathcal L/\partial S=-\lambda\left(1+\Delta\omega^N\right)C^N_S$,
\begin{align}
\dot p^S &= r\,p^S+\left(1+\Delta\omega^N\right)C^N_S,
\qquad\text{so}\qquad
\frac{\dot p^S}{p^S} = r+\frac{\left(1+\Delta\omega^N\right)C^N_S}{p^S}<r .
\label{eq:sp:A:costate:S}
\end{align}
This is the modified Hotelling rule: the in-situ rent grows more slowly than the interest rate, by
the marginal cost saving $-C^N_S>0$ that a larger known reserve confers on extraction, marked up by
the waste-base wedge on extraction effort.

\smalltitle{Discoveries.} With
$\partial\mathcal L/\partial X=-\lambda\left(1+\Delta\omega^D\right)C^D_X$,
\begin{align}
\dot p^X &= r\,p^X+\left(1+\Delta\omega^D\right)C^D_X,
\qquad\text{so}\qquad
p^X\left(t\right) = -\int_t^\infty \left(1+\Delta\omega^D\right)C^D_X\!\left(s\right)
e^{-\int_t^s r\,du}\,ds<0 .
\label{eq:sp:A:costate:X}
\end{align}
Cumulative discovery is a bad: each unit found today raises the cost of finding the next one forever
after, and $-p^X$ is the present value of that penalty.

\smalltitle{Pollution.} Write the marginal damage of the stock, in final-good units, as
\begin{align}
d &\equiv \frac{v'\!\left(P\right)}{\lambda}-\left(1-\Delta\omega^Y\right)F_P\;>\;0,
\label{eq:sp:A:damage}
\end{align}
the sum of the direct disutility and the output loss, the latter valued at the same wedge that
output carries everywhere else. Then
\begin{align}
\dot p^P &= \Big[r+\theta\!\left(P\right)+\theta'\!\left(P\right)P\Big]p^P-d,
\qquad\text{so}\qquad
p^P\left(t\right)=\int_t^\infty d\left(s\right)
e^{-\int_t^s\left[r+\theta+\theta'P\right]du}\,ds .
\label{eq:sp:A:costate:P}
\end{align}
The effective discount rate on damages is the interest rate plus the marginal decay rate
$\partial\left[\theta(P)P\right]/\partial P=\theta+\theta'P$. Because $\theta'<0$
by~\eqref{eq:sp:setup:pollution-motion}, sink saturation lowers that rate below the average decay rate
$\theta$, and $p^P$ is correspondingly larger than a naive $d/\left(r+\theta\right)$
calculation would suggest. Stability of the pollution block requires $\theta+\theta'P>0$; this is an
assumption on $\theta$, and it is maintained throughout.

\smalltitle{Embodied material.} With
$\partial\mathcal L/\partial M^K=\lambda\delta\left(p^M-p^W\right)$,
\begin{align}
\dot p^M &= \left(r+\delta\right)p^M-\delta p^W,
\qquad\text{so}\qquad
p^M\left(t\right)=\int_t^\infty \delta p^W\!\left(s\right)
e^{-\left(\int_t^s r\,du\right)-\delta\left(s-t\right)}ds .
\label{eq:sp:A:costate:M}
\end{align}
$p^M$ is the present value of the waste that the currently embodied stock will release as it
depreciates, discounted at $r+\delta$ because a fraction $\delta$ of the liability is discharged each
instant. It is positive whenever $p^W$ is, confirming the sign anticipated
in~\eqref{eq:sp:A:costates-def}, and it is strictly below the contemporaneous $p^W$ whenever $r>0$ ---
postponing a tonne of waste is worth something. Note that this equation is the one place in the
system where no wedge appears: the liability is denominated in tonnes throughout.

\smalltitle{The consumption Euler equation.} Combining~\eqref{eq:sp:A:foc:C}
and~\eqref{eq:sp:A:costate:K} and writing $\sigma\equiv-Cu''/u'$,
\begin{align}
\sigma\left(C\right)\frac{\dot C}{C}
&= \frac{\left(1-\Delta\omega^Y\right)F_K+\dot q}{q}-\delta-\rho
-\frac{d}{dt}\ln\!\left(1+\Delta\omega^C\right).
\label{eq:sp:A:euler}
\end{align}
The last term is what the failure of $\lambda=u'(C)$ costs. A wedge that is expected to grow --- which
is what a rising $p^W$ does --- tilts consumption toward the present, an effect with no counterpart in
Version B.

\smalltitle{Transversality.} The problem is closed by
\begin{align}
\lim_{t\to\infty}e^{-\rho t}\lambda q K &= 0,
&
\lim_{t\to\infty}e^{-\rho t}\lambda p^S S &= 0,
&
\lim_{t\to\infty}e^{-\rho t}\lambda p^X X &= 0,
\notag\\
\lim_{t\to\infty}e^{-\rho t}\lambda p^P P &= 0,
&
\lim_{t\to\infty}e^{-\rho t}\lambda p^M M^K &= 0 .
&&
\label{eq:sp:A:tvc}
\end{align}


\subsection{Reading the conditions}\label{sec:sp:A:reading}

\smalltitle{The materials arbitrage.} Everything in the materials block follows from one equation.
Rearranging the extraction condition~\eqref{eq:sp:A:foc:N}
\begin{align}
\underbrace{\left(1-\Delta\omega^Y\right)F_R-p^W+d^Wp^P+\zeta}_{\text{value of a secondary tonne}}
\;=\;
\underbrace{\left(1+\Delta\omega^N\right)C^N_N+p^S+\Omega^{N,S}p^W}_{\text{cost of a virgin tonne}}
\;+\;d^Wp^P .
\label{eq:sp:A:arbitrage}
\end{align}
A tonne of material can be obtained in two ways. Extracting it costs
$\left(1+\Delta\omega^N\right)C^N_N$ in resources, $p^S$ in foregone reserve value, and
$\Omega^{N,S}p^W$ in the waste cost of the overburden it drags up. Recovering it from waste instead
avoids all three, and additionally avoids the emission $d^W$ that the tonne would have caused had it
stayed untreated --- which is why $d^Wp^P$ appears on both sides and cancels out of the comparison of
\emph{sources} while remaining in the valuation of the recovered tonne.

The cancellation is the sharpest check available on the accounting. Note what has dropped out: $p^W$,
the shadow price of waste, appears on both sides of the arbitrage and is not what determines whether
recycling is worthwhile. It is worth being precise about why. A tonne of secondary material does not
\emph{remove} a tonne of waste from the economy; it postpones it. The tonne is recovered, used in
production, and returns as waste on the next pass. What recycling saves is not the waste cost but the
\emph{extraction} cost --- the resources, the rent, and the displaced overburden. This is the
substantive content of treating the materials balance as a ledger, and it survives the introduction of
the composition wedge unchanged.

Substituting~\eqref{eq:sp:A:arbitrage} back into the definition of $B$ gives the companion
statement for the production margin,
\begin{align}
\left(1-\Delta\omega^Y\right)F_R &= p^N+p^W-\zeta .
\label{eq:sp:A:FR}
\end{align}
The marginal product of materials in production, valued net of the waste base the resulting output
adds, covers the full cost of delivering a virgin tonne \emph{plus} the waste that tonne will
eventually become. Here $p^W$ does appear, and it must: a tonne drawn into the economy leaves as waste
exactly once, whatever route it takes on the way, and that is the point at which the waste cost is
charged.

\smalltitle{Are the wedges real?} \eqref{eq:sp:A:foc:C} puts a wedge between the marginal utility of
consumption and the shadow value of a unit of the final good, and~\eqref{eq:sp:A:foc:N} and its
companions do the same at every other point where the final good is spent. It is fair to ask whether
these are artefacts of the accounting. They are not, and the reason is that total mass is conserved
throughout: what consumption does under Version A is not reduce the mass drawn in, but change the
\emph{split} of a given mass between the durable stock and the current waste flow. That is a real
margin. The test is the multiplier: the wedge is proportional to $\chi\left(p^W-p^M\right)$, so it
vanishes when capital is not a sink ($\chi=0$ or $\phi^I=0$) and when postponement is worthless
($p^W=p^M$), which is exactly when the split it prices ceases to exist.
Section~\ref{sec:sp:B:nowedge} constructs a superficially similar wedge that \emph{is} an artefact, and
Section~\ref{sec:sp:B:compare} sets the two side by side; the comparison is the cleanest available
statement of what the process-level accounting buys.

\smalltitle{The signs of the prices.} Collecting~\eqref{eq:sp:A:costate:S}--\eqref{eq:sp:A:costate:M}:
$p^S>0$ (reserves are an asset), $p^X<0$ (cumulative discovery is a liability), $p^P>0$ (pollution is a
liability, so the costate $-\lambda p^P$ is negative as posited), and $p^M$ has the sign of $p^W$.
Only $p^W$ is of ambiguous sign. From~\eqref{eq:sp:A:foc:W-solved},
\begin{align}
p^W<0
\qquad\Longleftrightarrow\qquad
\varpi\alpha\Big[\left(1+\Delta\omega^N\right)C^N_N+p^S+d^Wp^P\Big]
\;>\;
\left(1+\Delta\omega^W\right)c^W\!\left(\varpi\right)+\left(1-\varpi+d^W\varpi\right)p^P,
\label{eq:sp:A:pW-sign}
\end{align}
that is, whenever the extraction a marginal tonne of waste displaces through recycling is worth more
than collecting, treating and emitting it costs. Since $p^S$ rises as reserves deplete
(Section~\ref{sec:sp:A:longrun}) while $c^W$ and $p^P$ need not, this is not a pathological case but a
natural late-transition one: waste becomes an asset.

Three objects flip sign together when it happens, and the coincidence is a useful internal check on
any numerical solution. By~\eqref{eq:sp:A:costate:M} the sign of $p^M$ follows the sign of $p^W$;
by~\eqref{eq:sp:A:foc:I} the wedge on capital then flips, so that $q>1$ --- embodying material in
capital becomes a cost rather than a saving; and by~\eqref{eq:sp:A:zeta} the composition wedge $\Delta$
flips with the same bracket, so that spending on consumption becomes socially cheap rather than dear.
All three share a single cause, and a solution in which they do not coincide is wrong. Any numerical
implementation must therefore leave $p^W$, $\Delta$ and every instrument built from them unbounded in
sign.

\smalltitle{The recycling block.} The two conditions~\eqref{eq:sp:A:foc:varpi}
and~\eqref{eq:sp:A:foc:KR} determine how circular the economy is, and they respond to a single
sufficient statistic, $B$. Condition~\eqref{eq:sp:A:foc:KR} allocates capital so that its
marginal product in production equals its marginal product in recycling, the latter being $a'(x)$
tonnes of secondary material each worth $B$; both sides are valued at the same wedge
$\left(1-\Delta\omega^Y\right)$ on the production side, so the wedge does not distort the allocation of
capital between the two uses, only the level at which it is worth holding.
Condition~\eqref{eq:sp:A:foc:varpi} pushes treatment to where its marginal cost equals the value of the
material it yields plus the emission it prevents. The two terms are the two distinct reasons to treat
waste, and they are separately identified --- an economy with $d^W=1$ (treatment that immobilises
nothing) still treats waste, purely to recycle it; an economy with $\alpha=0$ still treats waste,
purely to avoid emissions.

\begin{proposition}[Scarcity raises circularity, Version A]\label{prop:sp:A:circularity}
Hold $\left(1-\Delta\omega^Y\right)F_K$, $\left(1+\Delta\omega^W\right)$ and $p^P$ fixed and let
$B$ rise. Then $x$ rises, $a(x)$ rises, the marginal yield $\alpha$ rises, $\varpi$ rises,
and the recirculated share $a\varpi$ rises.
\end{proposition}

\begin{proof}
From~\eqref{eq:sp:A:foc:KR}, $a'(x)=\left(1-\Delta\omega^Y\right)F_K/B$ falls, and $a''<0$
gives $dx/dB>0$; $a'>0$ then gives $da/dB>0$. For the marginal yield,
$\alpha=a(x)-a'(x)x$ has $d\alpha/dx=a'-a''x-a'=-a''x>0$, so $\alpha$ rises with $x$ and hence with
$B$. Therefore $\alpha B$ rises, so the left-hand side
of~\eqref{eq:sp:A:foc:varpi} rises; $d^2c^T/d\varpi^2>0$ by~\eqref{eq:sp:setup:waste-cost} then gives
$d\varpi/dB>0$. Both factors of $a\varpi$ rise.
\end{proof}

The proof uses only the curvature of $a$ and the convexity of $c^T$, so the mechanism is untouched by
the composition wedge; what the wedge does is move the terms held fixed, which is a second-order
effect on the same margin. The mechanism of the paper is this proposition combined
with~\eqref{eq:sp:A:arbitrage} and~\eqref{eq:sp:A:costate:S}. Depletion raises the in-situ rent $p^S$;
\eqref{eq:sp:A:arbitrage} passes the rise one-for-one into $B$; and
Proposition~\ref{prop:sp:A:circularity} converts it into a higher recirculated share. By
\eqref{eq:sp:setup:ledger-collected} the circularity multiplier $\left(1-a\varpi\right)^{-1}$ rises
with it, so a given material throughput $R$ is supported by ever less virgin extraction. The
transition to a circular economy is, in this model, a scarcity-driven reallocation of capital, not a
pollution-driven one --- the pollution channel enters $B$ only through $d^Wp^P$, which is
bounded by the damage function, whereas $p^S$ is not bounded at all. Under Version A it is, in
addition and at second order, composition-driven.


\subsection{Stationarity and the long run}\label{sec:sp:A:longrun}

It is natural to look for a steady state. The model as specified does not have one, and the reason is
instructive.

\begin{proposition}[No interior stationary point, Version A]\label{prop:sp:A:nosteadystate}
At any point where all five states of Problem~\ref{prob:sp:A} are stationary, $D=N=0$ and
$W\left(1-a\varpi\right)=0$. Since $a(x)<1$ for all finite $x$, it follows that $W=R=R^R=0$: the
economy uses no materials.
\end{proposition}

\begin{proof}
$\dot X=D=0$ forces $D=0$, and then $\dot S=D-N=0$ forces $N=0$, and hence $\mathcal N=0$. $\dot K=0$
gives $I=0$ and $G=\delta K$; $\dot M^K=0$ then gives $\Omega\phi^IG=\delta M^K$, so the two capital
terms cancel in the ledger~\eqref{eq:sp:A:ledger}, leaving $W=R$. But $R=N+a\varpi W=a\varpi W$, so
$W\left(1-a\varpi\right)=0$. With $\varpi\le1$ and $a<1$ the factor is strictly positive.
\end{proof}

The culprit can be identified exactly, and it is not the closure: the proof uses only the state
equations and the ledger, both of which are common to the two versions. A stationary reserve with
positive extraction requires $D=N>0$, since $\dot S=D-N$; extraction must be replenished by discovery.
But $\dot X=D$ makes cumulative discovery a monotone state, so $D>0$ is itself incompatible with
stationarity. The only route to an interior rest point is therefore closed off by the single
assumption that exploration costs depend on cumulative past discovery. Drop it --- set $C^D_X=0$, so
that $X$ ceases to be a state --- and a genuine steady state with $D=N>0$ exists; this is
switch~\ref{sw:sp:CDX} in Section~\ref{sec:sp:schemes}, and it is the cheapest available repair. Note
that removing exploration altogether does \emph{not} repair it: with $\dot S=-N$, stationarity forces
$N=0$ directly. Section~\ref{sec:sp:setup:primitives} already flagged the corresponding structural
feature --- there is no upper bound on $S$, so the resource is never exhausted, only ever more costly
--- and Proposition~\ref{prop:sp:A:nosteadystate} is its dynamic counterpart.

That is in fact the substantive content. Read the proposition in the limit rather than at a point. As
$t\to\infty$ the rent $p^S$ grows, $B$ grows with it by~\eqref{eq:sp:A:arbitrage}, and
Proposition~\ref{prop:sp:A:circularity} sends $x\to\infty$, $a\to1$ and $\varpi\to1$. Along such a path
$N\to0$ while $R$ need not: from~\eqref{eq:sp:setup:ledger-collected} with $\dot M^K\to0$ and
$\mathcal N\to0$, $R\left(1-a\varpi\right)\to0$ with $R$ itself indeterminate, its value being whatever
material throughput the recycling sector can sustain. The emission flow~\eqref{eq:sp:A:Xi} tends to
$\left(1-\varpi+d^W\varpi\right)W-d^Wa\varpi W\to0$, so the limiting economy is materially closed and
emission-free, and it operates at positive scale. This is the sense in which a fully circular economy
is the model's long run: not a steady state that is reached, but a boundary that is approached, and
approached only because extraction becomes infinitely expensive.

Two further remarks on the limit are specific to Version A. As $a\varpi\to1$ and $N\to0$ with $R$
bounded, $\chi=\Phi^IG/R$ need not vanish, so $\Delta$ need not vanish either: the composition wedge
survives into the limit, and the two versions do not converge as the economy becomes circular. And
$\Omega=R/\left(\mathcal D+\phi^IG\right)$ remains strictly positive along the path, so the accounting
never degenerates.

\smalltitle{A caveat on the multiplier.} With $\dot M^K=0$ the collected
ledger~\eqref{eq:sp:setup:ledger-collected} reads $N+\mathcal N=W\left(1-a\varpi\right)$, which makes
the quantitative stakes of the circularity multiplier visible. As $a\varpi\to1$ a given $N$ supports an
unbounded $W$, so the model's mass flows become arbitrarily sensitive to the recirculation assumption in
exactly the region the paper is about. The assumption in question is that recycled material re-enters
production contemporaneously --- $R^R=a\varpi W$ at the same instant --- which in continuous time is
innocuous but in a discrete-time implementation ties the number of circuits per period to the period
length. It is flagged here because $R/N$ and $W/N$ should be reported alongside any headline result,
and their sensitivity to the period length checked; see~\ref{op:sp:period}.


\subsection{The complete system}\label{sec:sp:A:system}

For reference, and as the specification the quantitative implementation must reproduce, the
optimality system of Version A consists of:
\begin{itemize}[leftmargin=1.6em,itemsep=0.15em]
\item five state equations~\eqref{eq:sp:A:states-motion} in $K,S,X,P,M^K$;
\item six price equations: \eqref{eq:sp:A:costate:K} for $q$, \eqref{eq:sp:A:costate:S} for $p^S$,
\eqref{eq:sp:A:costate:X} for $p^X$, \eqref{eq:sp:A:costate:P} for $p^P$, \eqref{eq:sp:A:costate:M}
for $p^M$, and \eqref{eq:sp:A:euler} for $C$ (equivalently for $\lambda$);
\item five static conditions determining $I,N,D,\varpi,K^R$: \eqref{eq:sp:A:foc:I},
\eqref{eq:sp:A:foc:N}, \eqref{eq:sp:A:foc:D}, \eqref{eq:sp:A:foc:varpi}, \eqref{eq:sp:A:foc:KR};
\item three definitions of prices: \eqref{eq:sp:A:foc:W} for $p^W$, \eqref{eq:sp:A:zeta} for $\zeta$
and $\Delta$, and~\eqref{eq:sp:A:damage} for $d$;
\item the goods constraint~\eqref{eq:sp:A:goods}, the ledger~\eqref{eq:sp:A:ledger}, and the intensity
definition~\eqref{eq:sp:A:intensity}, which must be solved \emph{jointly} for $W$ and $\Omega$ ---
equivalently, $W$ is the positive root of the quadratic~\eqref{eq:sp:setup:ledger-quadratic} and
$\Omega$ follows from~\eqref{eq:sp:setup:Omega-phi};
\item the definitions~\eqref{eq:sp:A:definitions} for $R^R,R,Y,G$ and~\eqref{eq:sp:A:Xi} for $\Xi$;
\item initial conditions on the five states and the transversality conditions~\eqref{eq:sp:A:tvc}.
\end{itemize}
There is no complementary-slackness block: the mass restriction $\Phi^IG\le R$ is vacuous here, by~(i)
of Section~\ref{sec:sp:A:closure}. In exchange, $\Omega$ is a genuine unknown of the system rather than
a quantity computed after the fact, and it enters the optimality conditions through $\Delta$ at every
point where the final good is spent. This is the structural difference from Version B, where the
accounting objects do not feed back and can be recovered afterwards.

\begin{remark}[The reduced formulation]\label{rem:sp:A:reduced}
Substituting the positive root of~\eqref{eq:sp:setup:ledger-quadratic} and the
solution~\eqref{eq:sp:setup:Omega-phi} for $\Omega$ into the goods constraint, and dropping $W$ and
$\Omega$ from the control set, gives a problem in six controls and the same five states, with the same
solution. The two formulations are equivalent because \eqref{eq:sp:A:foc:W} and~\eqref{eq:sp:A:foc:Omega}
are precisely the conditions that make the substitution valid; what is lost is the two prices
themselves. Since $p^W$ is the object the market economy implements through a waste tax, and since
$\zeta$ is what the entire difference between the two versions is expressed in terms of, both are
worth keeping.
\end{remark}

\begin{remark}[What is not verified]\label{rem:sp:A:soc}
The conditions above are necessary. Sufficiency requires concavity of the maximised Hamiltonian in the
states, which is not established here: $F$ is concave and $u,-v$ are, but $\mathcal R$ enters through a
constant-returns function of two controls, the ledger is bilinear in $\varpi$ and $W$, and the intensity
definition~\eqref{eq:sp:A:intensity} is bilinear in $\Omega$ and the flows in $\mathcal D$. Uniqueness of
the optimum is likewise not established, and it is invoked later when the market implementation is
claimed to reproduce \emph{the} planner's allocation. Both are open; see~\ref{op:sp:soc}.
\end{remark}
